\section{Preliminaries}

\subsection{Radial Basis Function (RBF) Neural Networks}
\label{subsec:rbf_nn}

In this paper, Radial Basis Function (RBF) neural networks (NNs) are employed to approximate the unknown nonlinear lumped uncertainty in the robotic manipulator dynamics.

For an $n$-degree-of-freedom (DOF) robotic system, the input vector of the RBF neural network is selected as the tracking error state:

\begin{equation}
	x
	=
	X
	=
	\begin{bmatrix}
		e\\
		\dot e
	\end{bmatrix}
	\in\mathbb{R}^{2n},
	\label{eq:rbf_input}
\end{equation}

where $e=q-q_d$ and $\dot e=\dot q-\dot q_d$ denote the position tracking error and velocity tracking error, respectively.

An RBF neural network consists of a nonlinear mapping from the input vector $x$ to the hidden layer and a linear combination of the hidden-layer outputs. The estimated lumped uncertainty is expressed as:

\begin{equation}
	\hat{f}(x)=\hat{W}^{T}h(x),
	\label{eq:rbf_estimation}
\end{equation}

where $\hat{W}\in\mathbb{R}^{m\times n}$ is the estimated weight matrix, $m$ represents the number of hidden neurons, and

\begin{equation}
	h(x)
	=
	[h_1(x),h_2(x),\dots,h_m(x)]^T
	\in\mathbb{R}^{m}
\end{equation}

denotes the hidden-layer activation vector.

The Gaussian radial basis function of the $i$-th hidden neuron is defined as:

\begin{equation}
	h_i(x)
	=
	\exp
	\left(
	-\frac{\|x-c_i\|^2}{2b_i^2}
	\right),
	\quad i=1,2,\dots,m,
\end{equation}

where $c_i\in\mathbb{R}^{2n}$ represents the center vector of the $i$-th Gaussian basis function and $b_i>0$ denotes the corresponding width parameter.

According to the universal approximation theorem, for any continuous nonlinear function $f(x)$ defined on a compact set $\Omega_x\subset\mathbb{R}^{2n}$, there exists an ideal constant weight matrix $W^*\in\mathbb{R}^{m\times n}$ such that:

\begin{equation}
	f(x)
	=
	{W^*}^{T}h(x)
	+
	\epsilon,
	\label{eq:rbf_approximation}
\end{equation}

where $\epsilon\in\mathbb{R}^{n}$ denotes the inherent approximation error of the RBF neural network.

The ideal weight matrix is defined as:

\begin{equation}
	W^*
	\triangleq
	\arg\min_W
	\left[
	\sup_{x\in\Omega_x}
	\|f(x)-W^Th(x)\|
	\right].
\end{equation}

Furthermore, the approximation error is assumed to be bounded on the compact set $\Omega_x$:

\begin{equation}
	\|\epsilon\|\leq\epsilon_N,
\end{equation}

where $\epsilon_N>0$ is an unknown positive constant.

\begin{remark}
	The linearly parameterized structure of the RBF neural network indicates that the network output is linear with respect to the adjustable weight matrix $\hat{W}$, while the nonlinear approximation capability is provided by the basis function vector $h(x)$. This property facilitates the design of adaptive weight update laws based on Lyapunov stability analysis.
\end{remark}


\subsection{Notation}

The dimensions of the main variables used throughout this paper are summarized in Table~\ref{tab:notation}.

\begin{table}[!t]
	\caption{Dimensions of the Main Variables}
	\label{tab:notation}
	\centering
	\renewcommand{\arraystretch}{1.15}
	\begin{tabular}{ll}
		\hline
		\textbf{Symbol} & \textbf{Dimension}\\
		\hline
		$n$
		& Scalar
		\\
		$q,\dot q,\ddot q$
		& $\mathbb{R}^{n\times1}$
		\\
		$q_d,\dot q_d,\ddot q_d$
		& $\mathbb{R}^{n\times1}$
		\\
		$e,\dot e$
		& $\mathbb{R}^{n\times1}$
		\\
		$x=X=
		\begin{bmatrix}
			e\\
			\dot e
		\end{bmatrix}$
		& $\mathbb{R}^{2n\times1}$
		\\
		$m$
		& Scalar
		\\
		$h(x)$
		& $\mathbb{R}^{m\times1}$
		\\
		$c_i$
		& $\mathbb{R}^{2n\times1}$
		\\
		$C=[c_1,c_2,\dots,c_m]$
		& $\mathbb{R}^{2n\times m}$
		\\
		$b_i$
		& Scalar
		\\
		$W^*,\hat W,\tilde W$
		& $\mathbb{R}^{m\times n}$
		\\
		$f(x),\hat f(x)$
		& $\mathbb{R}^{n\times1}$
		\\
		$\epsilon$
		& $\mathbb{R}^{n\times1}$
		\\
		$P$
		& $\mathbb{R}^{2n\times2n}$
		\\
		$B=
		\begin{bmatrix}
			0\\
			M_0^{-1}(q)
		\end{bmatrix}$
		& $\mathbb{R}^{2n\times n}$
		\\
		$d(t)$
		& $\mathbb{R}^{n\times1}$
		\\
		\hline
	\end{tabular}
\end{table}